\documentclass{beamer}

\usepackage{cuhkbeamer}


\begin{document}

% ====================== Slide 1: TITLE SLIDE ======================
\begin{frame}
    \title{Introduction to LAE}
    \author{Chiu Yu KO}
    \institute{Low Altitude Economics}
    \date{Week 1}
    
  \titlepage
\end{frame}

% ====================== Slide 2: FORMAT & OBJECTIVES ======================
\begin{frame}[t]
  \frametitle{Learning Objectives}
  By the end of this class, you should be able to:
  \begin{itemize}
    \item \textbf{Define} the Low Altitude Economy (LAE) in plain English.
    \item \textbf{Identify} the main business opportunities within the LAE.
    \item \textbf{Map} who creates value and who gets paid in the LAE ecosystem.
    \item \textbf{Explain} how Hong Kong's regulatory sandbox can support market development.
    \item \textbf{Use} simple calculations to assess whether an LAE service might be commercially viable.
  \end{itemize}
\end{frame}

% ====================== Slide 3: ICEBREAKER (REVISED WITH FACTS) ======================
\begin{frame}[t]
  \frametitle{Why LAE Matters for Hong Kong}
  \textbf{Imagine two possible services in Hong Kong:}
  \begin{itemize}
      \item A drone carries urgent medical supplies across Victoria Harbour, avoiding road congestion.
      \item An electric air taxi connects a business district with the airport, reducing door-to-door travel time.
  \end{itemize}
  \vspace{0.3cm}
  \textbf{What could this mean for the economy?}
  \begin{itemize}
    \item Low-altitude airspace adds a \textbf{third dimension} to the transport network.
    \item One 2025 study projects that Hong Kong's LAE-related industries could generate approximately \textbf{HK\$350 billion in annual value added by 2050}, with supporting services contributing a substantial share.
    \item \textbf{Quick poll (2 min):} ``Where is Hong Kong most likely to create value in the low-altitude economy?''
  \end{itemize}
  \vfill
  {\tiny Source: Our Hong Kong Foundation, \textit{Taking Flight: Forging a Future for Hong Kong's Low-Altitude Economy} (2025).}
\end{frame}

% ====================== Week 1 - Session 1 Title ======================
\begin{frame}[c]
  \begin{center}
    \huge\color{cuhkpurple}\textbf{Session 1}\\
    \vspace{0.5cm}
    \Large Foundations, Definitions \& Economic Significance
  \end{center}
\end{frame}

% ====================== Slide 3.5: WHAT IS LAE (VISUAL + FACT) ======================
\begin{frame}[t]
  \frametitle{What is the Low Altitude Economy?}
  \textbf{LAE = Economic activity that uses low-altitude airspace, aircraft, infrastructure, and related services.}
  \begin{columns}[T, onlytextwidth]
    \begin{column}{0.35\textwidth}
      \raggedleft
      \begin{tikzpicture}[scale=0.75, every node/.style={font=\small}]
        \useasboundingbox (-2.2,-1.2) rectangle (2.2,4);
        \fill[gray!10] (-2.2,0) rectangle (2.2,4);
        \fill[gray!30] (-2.2,-0.5) rectangle (2.2,0);
        \fill[cuhkpurple!20] (-2.0,0) rectangle (-1.2,2.5);
        \fill[cuhkpurple!20] (-0.5,0) rectangle (0.4,3.0);
        \fill[cuhkpurple!20] (1.1,0) rectangle (2.0,2.0);
        \draw[thick, cuhkpurple, <->] (-2.2,3.5) -- (2.2,3.5) node[midway, fill=white, inner sep=1pt] {\tiny typically below 1,000 m};
        \node[draw, circle, minimum size=4pt, fill=black] (drone) at (-0.8,1.6) {};
        \draw[thick] (-1.0,1.6) -- (-0.6,1.6);
        \draw[thick] (-0.9,1.8) -- (-0.7,1.8);
        \node[left=1pt of drone, font=\tiny\bfseries] {Drone};
        \fill[cuhkgold] (1.1,2.0) rectangle (2.0,2.15);
        \node[align=center, font=\tiny\bfseries] at (1.55, 2.3) {Vertiport};
        \node[align=center, text width=3.2cm, font=\tiny] at (0,-0.9) {\textbf{eVTOL}: electric vertical take-off \\ and landing aircraft};
      \end{tikzpicture}
    \end{column}
    \begin{column}{0.65\textwidth}
      \hspace{-0.5cm}
      \begin{itemize}
        \item \textbf{Aircraft}: Drones, helicopters, and eVTOL aircraft.
        \item \textbf{Physical infrastructure}: Vertiports, charging facilities, and maintenance hubs.
        \item \textbf{Digital infrastructure}: Communications, traffic-management software, and data platforms.
        \item \textbf{Applications and services}: Logistics, inspection, emergency response, passenger transport, finance, and insurance.
      \end{itemize}
      \vspace{0.25cm}
      \small\textit{The exact altitude boundary varies across policy and industry contexts.}
    \end{column}
  \end{columns}
\end{frame}

% ====================== REVISED SLIDE 7: LAE vs TRADITIONAL AVIATION ======================
\begin{frame}[t]
  \frametitle{How LAE Differs from Traditional Aviation}
 
  \textbf{Recall:} LAE covers economic activity using low-altitude airspace and its supporting ecosystem.
 
  Traditional aviation (e.g., Cathay Pacific) operates very differently:
  \begin{itemize}
    \item Uses massive, expensive aircraft
    \item Departs from giant, centralized airports
    \item Flies long distances at high altitude
  \end{itemize}
 
  \vspace{0.3cm}
  \textbf{The LAE Shift: Key Business Implications}
  \begin{itemize}
    \item \textbf{Short-distance, high-frequency operations} inside or between cities
    \item \textbf{Entire ecosystem} required: charging pads, digital traffic software, insurance, and maintenance
    \item \textbf{Analogy:} Think of LAE less like a traditional airline and more like an Uber or broadband internet network for the sky
  \end{itemize}
 
  \textit{Result: LAE services have different cost structures, operating models, and regulatory needs.}
\end{frame}

% ====================== Slide 8: CLARIFICATION ======================
\begin{frame}[t]
  \frametitle{Jargon Buster: LAE, UAM, and AAM}

  You will hear these acronyms constantly. Here is what they actually mean:

  \begin{itemize}
    \item \textbf{LAE (Low Altitude Economy):} The entire umbrella term we defined earlier. Includes drones, flying taxis, infrastructure, and the software running it all.
    \item \textbf{UAM (Urban Air Mobility):} A specific subset of LAE. This refers \textit{only} to moving people or cargo over short distances inside a crowded city (like Central to Kowloon).
    \item \textbf{AAM (Advanced Air Mobility):} A broader term for using new aircraft to connect cities to suburbs or rural areas.
    \item \textbf{Course focus:} LAE economics starting \textbf{HK-first}, then expanding to the \textbf{Greater Bay Area (GBA)} and \textbf{the rest of China}.
  \end{itemize}
\end{frame}

% ====================== Slide 9: CORE USE CASES (REVISED WITH FACTS) ======================
\begin{frame}[t]
  \frametitle{Four Ways to Create Value in the Sky}
  LAE services must solve a customer problem better than existing alternatives.
  \vspace{0.25cm}
  \begin{enumerate}
    \item \textbf{Cargo and logistics:} High-priority items for which speed, reliability, or access matters.
    \item \textbf{Passenger air mobility:} Premium short-distance travel where door-to-door time savings justify the fare.
    \item \textbf{Industrial inspection:} Inspection of bridges, buildings, power infrastructure, and remote assets with less disruption and risk.
    \item \textbf{Emergency and public services:} Medical delivery, emergency response, search and rescue, and environmental monitoring.
  \end{enumerate}
  \vspace{0.35cm}
  \textbf{Managerial question:} Which applications create enough value to cover the cost, risk, and regulatory requirements?
\end{frame}



% ====================== Slide 4: CONCEPT PRIMER 1 ======================
\begin{frame}[t]
  \frametitle{Concept Primer: ``Value of Time'' \& ``NPV''}
  \textbf{1. Value of Time (VOT)}
  \begin{itemize}
    \item VOT expresses how much a customer is willing to pay to save travel time.
    \item It varies across customers, purposes, and situations. Saving 20 minutes may be worth much more for an urgent medical delivery than for a routine trip.
  \end{itemize}
  \vspace{0.25cm}
  \textbf{2. Net Present Value (NPV)}
  \begin{itemize}
    \item NPV asks whether the future cash generated by a project is sufficient to justify the investment made today.
    \item A positive NPV suggests that the project may be financially attractive. A negative NPV means that it does not justify the investment under the assumptions used.
  \end{itemize}
\end{frame}

% ====================== REVISED SLIDE 5: NPV FORMULA (NOW CONSISTENT) ======================
\begin{frame}[t]
  \frametitle{Applied Tool: Assessing Market Viability}
  \begin{block}{Step 1: Does the service create value for customers?}
  \[
    \text{Value of Time Saved per Trip}=VOT\times\Delta T
  \]
  \end{block}
  \begin{block}{Step 2: Can the operator capture enough of that value?}
  \[
    \text{Annual Operating Profit}=\text{Annual Revenue}-\text{Annual Operating Cost}
  \]
  \end{block}
  \begin{block}{Step 3: Does the project justify the initial investment?}
  \[
    \text{NPV}=-\text{Initial Investment}+\text{Present Value of Future Profits}
  \]
  \end{block}
  \textbf{Core lesson:} Time savings support willingness to pay, but fare, demand, costs, and investment determine financial viability.
\end{frame}

% ====================== REVISED SLIDE 6: BACK-OF-ENVELOPE (NPV) – NOW CONSISTENT ======================
\begin{frame}[t, shrink=15]
  \frametitle{Back-of-the-Envelope: How Friction Destroys Value}
  \textbf{Scenario:} A proposed air-taxi service connects HKIA and Central. The customer's time is worth HK\$1,000 per hour.
  \begin{columns}[T]
    \begin{column}{0.48\textwidth}
      \begin{block}{Case A: Smooth Journey}
      \begin{itemize}
        \item Flight time: 10 minutes
        \item Boarding and waiting: 5 minutes
        \item Total journey: 15 minutes
        \item Time saved: 30 minutes = 0.5 hours
        \item Value of time saved: \(1000\times0.5=\text{HK\$500}\)
      \end{itemize}
      \end{block}
    \end{column}
    \begin{column}{0.48\textwidth}
      \begin{block}{Case B: Friction Shock}
      \begin{itemize}
        \item Flight time: 10 minutes
        \item Boarding and waiting: 25 minutes
        \item Total journey: 35 minutes
        \item Time saved: 10 minutes \(\approx0.167\) hours
        \item Value of time saved: \(1000\times0.167\approx\text{HK\$167}\)
      \end{itemize}
      \end{block}
    \end{column}
  \end{columns}
  \vspace{0.25cm}
  \textit{Takeaway: A fast aircraft does not guarantee a valuable service. Waiting, transfers, and access time can sharply reduce the fare premium that customers will accept.}
  \vfill
  {\tiny Illustrative assumptions for classroom analysis; not current market data.}
\end{frame}

% ====================== Slide 10: CONCEPT PRIMER 2 ======================
\begin{frame}[t]
  \frametitle{Concept Primer: How Customers Compare Transport Options}
  When choosing between the MTR, a taxi, or an air taxi, customers consider more than the ticket price.
  \begin{itemize}
    \item \textbf{Money cost}: Fare and related charges.
    \item \textbf{Time cost}: In-vehicle time, waiting, boarding, and transfers.
    \item \textbf{Other factors}: Reliability, convenience, comfort, and perceived safety.
  \end{itemize}
  \vspace{0.35cm}
  Economists combine money and time into a simple measure called \textbf{generalized cost}.
\end{frame}

% ====================== Slide 11: STANDALONE FORMULA ======================
\begin{frame}[t]
  \frametitle{Applied Tool: Generalized Cost}
  \begin{block}{The Formula}
  \[
    \text{Generalized Cost}=\text{Ticket Price}+VOT\times\text{Travel Time}
  \]
  \end{block}
  \begin{itemize}
    \item The formula converts travel time into a monetary value.
    \item All else equal, the customer prefers the option with the lower generalized cost.
    \item A faster mode can charge more, but only up to the value of the time it saves.
  \end{itemize}
  \begin{block}{Managerial Lesson}
    Route planners must consider the complete journey, not only the time spent in the aircraft.
  \end{block}
\end{frame}

% ====================== Slide 12: BACK-OF-ENVELOPE (UTILITY) ======================
\begin{frame}[t, shrink=10]
  \frametitle{Back-of-the-Envelope: Finding the Maximum Ticket Price}
  \textbf{Scenario:} What is the most a time-sensitive traveller would pay for an air taxi before choosing the Airport Express instead?
  \textbf{Assumption:} VOT = HK\$10 per minute.
  \begin{columns}[T]
    \begin{column}{0.48\textwidth}
      \begin{block}{Airport Express}
      \begin{itemize}
        \item Ticket price: HK\$115
        \item Journey time: 24 minutes
        \item \(GC_{MTR}=115+10(24)\)
        \item \textbf{Generalized cost = HK\$355}
      \end{itemize}
      \end{block}
    \end{column}
    \begin{column}{0.48\textwidth}
      \begin{block}{Air Taxi}
      \begin{itemize}
        \item Journey time: 8 minutes
        \item Match the MTR generalized cost:
        \item \(P+10(8)=355\)
        \item \textbf{Maximum ticket price = HK\$275}
      \end{itemize}
      \end{block}
    \end{column}
  \end{columns}
  \vspace{0.25cm}
  \textit{Takeaway: Under these simplified assumptions, HK\$275 makes the traveller indifferent. Actual demand also depends on waiting, access, reliability, comfort, and perceived safety.}
  \vfill
  {\tiny Illustrative assumptions for classroom analysis; not current market data.}
\end{frame}

% ====================== Slide 13: ECONOMIC SCALE HK ======================
\begin{frame}[t]
  \frametitle{Hong Kong's Possible Role: The ``Service Layer''}
  \textbf{If LAE creates substantial value, where might Hong Kong hold an advantage?}
  \begin{itemize}
    \item The GBA has significant strengths in aircraft manufacturing, electronics, supply chains, and large-scale testing.
    \item Hong Kong may contribute high-value supporting services that help LAE businesses operate, finance growth, and expand internationally:
    \begin{itemize}
      \item \textbf{Finance and leasing}
      \item \textbf{Insurance and risk management}
      \item \textbf{Legal, certification, and professional services}
      \item \textbf{Digital platforms and international business development}
    \end{itemize}
  \end{itemize}
  \textit{This is a strategic proposition to evaluate, not a predetermined outcome.}
\end{frame}

% ====================== Slide 14: SESSION 2 TITLE ======================
\begin{frame}[c]
  \begin{center}
    \huge\color{cuhkpurple}\textbf{Session 2}\\
    \vspace{0.5cm}
    \Large Who Makes What? (The Value Chain)
  \end{center}
\end{frame}

% ====================== Slide 15: JARGON BUSTER ======================
\begin{frame}[t]
  \frametitle{Jargon Buster: The LAE Ecosystem}
  \begin{itemize}
    \item \textbf{OEM (Original Equipment Manufacturer):} A company that produces the aircraft or key components.
    \item \textbf{Vertiport:} A site designed for eVTOL take-off, landing, passenger handling, charging, and related operations.
    \item \textbf{UTM / U-space:} Digital services and procedures that support the safe and efficient management of unmanned-aircraft operations.
    \item \textbf{Operator:} The organization that provides the transport, logistics, inspection, or public service to the customer.
  \end{itemize}
\end{frame}



% ====================== REVISED SLIDE 16: LAE VALUE CHAIN (VISUAL) ======================



\begin{frame}[t]
  \frametitle{The LAE Ecosystem -- Who Gets Paid?}
  \begin{center}
  \resizebox{0.92\linewidth}{!}{%
  \begin{tikzpicture}[every node/.style={font=\small, align=center}]
    \tikzset{box/.style={rectangle, rounded corners=5pt, minimum width=3.2cm, minimum height=0.8cm, draw=cuhkpurple, fill=cuhklight, thick}}
    
    \node[box] (oem) at (-3.8, 1.8) {Aircraft and\\Components};
    \node[box] (infra) at (0, 1.8) {Vertiports, Energy\\and Maintenance};
    \node[box] (digital) at (3.8, 1.8) {UTM, Communications\\and Data};
    \node[box, fill=cuhkgold!25] (operator) at (0, 0.4) {Operators};
    \node[box] (users) at (0, -0.9) {Customers and Public Agencies};
    
    % Scaled down width and broken into two readable lines
    \node[box, text width=9.5cm] (services) at (0, -2.1) {Supporting Services: Finance, Leasing, Insurance,\\Certification, Legal and Professional Services};

    \draw[->, thick, cuhkpurple] (oem) -- (operator);
    \draw[->, thick, cuhkpurple] (infra) -- (operator);
    \draw[->, thick, cuhkpurple] (digital) -- (operator);
    \draw[<->, thick, cuhkgold] (operator) -- (users);
    \draw[->, thick, cuhkpurple] (services) -- (operator);
    \draw[->, thick, cuhkpurple] (services.west) to[bend left=20] (oem.south);
    \draw[->, thick, cuhkpurple] (services.east) to[bend right=20] (digital.south);
  \end{tikzpicture}%
  }
  \end{center}
  \vfill
  \centering\small\textit{Managerial question: Which participants create value, bear risk, and have bargaining power?}
\end{frame}

% ====================== Slide 16.1: APPLIED TOOL (VALUE CAPTURE) ======================
\begin{frame}[t, shrink=10]
  \frametitle{Applied Tool: How the Ticket Price is Shared}
  \begin{block}{A Simple Price Decomposition}
  \[
  \begin{split}
  \text{Ticket Price}={}&\text{Aircraft Cost}+\text{Vertiport Fee}+\text{Energy}\\
  &+\text{UTM Fee}+\text{Maintenance and Labour}\\
  &+\text{Insurance and Overhead}+\text{Operator Margin}
  \end{split}
  \]
  \end{block}
  \begin{itemize}
    \item Different participants may charge usage fees, lease payments, or service fees.
    \item The operator faces direct demand risk because it must attract customers and fill available capacity.
    \item Infrastructure and platform providers may face less direct passenger risk, but they still depend on traffic volume and bear investment, technology, and regulatory risks.
  \end{itemize}
  \textbf{Core lesson:} Creating value does not guarantee capturing value. Contract design and bargaining power determine how revenue is shared.
\end{frame}

% ====================== Slide 16.2: BACK-OF-ENVELOPE (PIE CHART) ======================
\begin{frame}[t, shrink=10]
  \frametitle{Illustrative Unit Economics: Where Might HK\$300 Go?}
  \textbf{Scenario:} A passenger pays HK\$300 for an air-taxi trip. Assume the fare is allocated as follows.
  \begin{columns}[T]
    \begin{column}{0.48\textwidth}
      \begin{block}{Payments to Partners}
      \begin{itemize}
        \item Vertiport fee: HK\$100
        \item UTM and digital services: HK\$20
        \item Insurance and financing: HK\$30
        \item Aircraft lease allocation: HK\$50
        \item \textbf{Total partner payments: HK\$200}
      \end{itemize}
      \end{block}
    \end{column}
    \begin{column}{0.48\textwidth}
      \begin{block}{Operator Economics}
      \begin{itemize}
        \item Cash remaining: HK\$100
        \item Energy: HK\$20
        \item Maintenance and labour: HK\$50
        \item \textbf{Operator margin: HK\$30}
      \end{itemize}
      \end{block}
      \textit{If all costs remain unchanged, cutting the fare by HK\$50 turns a HK\$30 margin into a HK\$20 loss.}
    \end{column}
  \end{columns}
  \vfill
  {\tiny Hypothetical figures for classroom discussion only. Actual costs and contracts may differ substantially.}
\end{frame}

% ====================== Slide 17: INCENTIVES & FLOWS ======================
\begin{frame}[t]
  \frametitle{Stakeholder Conflict: Different Objectives}
  An \textbf{externality} arises when an activity imposes costs or benefits on others that are not fully reflected in the decision-maker's private costs or revenues.
  \begin{itemize}
    \item \textbf{Operators:} Want enough flights and passengers to cover fixed costs.
    \item \textbf{Investors:} Want a credible route to approval, revenue, and scale.
    \item \textbf{Government:} Balances innovation, safety, efficiency, and public interest.
    \item \textbf{Communities:} Care about noise, privacy, visual impact, and fair access to benefits.
  \end{itemize}
  \textbf{Managerial challenge:} How can a service reach commercial scale while maintaining public acceptance?
\end{frame}

% ====================== Slide 17.1: BACK-OF-ENVELOPE (CURFEW) ======================
\begin{frame}[t, shrink=10]
  \frametitle{Back-of-the-Envelope: How Operating Restrictions Affect Cost}
  \textbf{Scenario:} Noise concerns lead to a restriction on evening and overnight flights. Assume a fixed aircraft lease cost of HK\$10,000 per day.
  \begin{columns}[T]
    \begin{column}{0.48\textwidth}
      \begin{block}{Case A: Wider Operating Window}
      \begin{itemize}
        \item Flights per day: 50
        \item Fixed cost per flight:
        \item HK\$10,000 / 50 = \textbf{HK\$200}
      \end{itemize}
      \end{block}
    \end{column}
    \begin{column}{0.48\textwidth}
      \begin{block}{Case B: Restricted Operating Window}
      \begin{itemize}
        \item Flights per day: 20
        \item Fixed cost per flight:
        \item HK\$10,000 / 20 = \textbf{HK\$500}
      \end{itemize}
      \end{block}
    \end{column}
  \end{columns}
  \vspace{0.25cm}
  \textit{Takeaway: Public acceptance is not merely a public-relations issue. Operating restrictions can directly affect utilization, unit costs, and commercial viability.}
  \vfill
  {\tiny Illustrative assumptions for classroom analysis; not current policy or market data.}
\end{frame}

% ====================== Slide 18: SESSION 3 TITLE ======================
\begin{frame}[c]
  \begin{center}
    \huge\color{cuhkpurple}\textbf{Session 3}\\
    \vspace{0.5cm}
    \Large Hong Kong Policy, "The Sandbox", and Technology
  \end{center}
\end{frame}

% ====================== Slide 19: THE AVIATION PROBLEM ======================
\begin{frame}[t]
  \frametitle{Concept Primer: The Aviation Safety Problem}
  
  \textbf{Why don't we have flying taxis everywhere already?}

  \vspace{0.3cm}
  \begin{itemize}
      \item Aviation is the safest form of travel in the world because of incredibly strict regulations (often called "Part 25" certification). 
      \item To get a new airplane certified for the public, a company must spend \textbf{up to 10 years} and billions of dollars proving it won't crash.
      \item \textbf{The Problem:} Startups building drones do not have billions of dollars. If they have to wait 10 years to fly their first paying customer, they will run out of money and die.
  \end{itemize}
\end{frame}

% ====================== Slide 20: POLICY MOMENTUM ======================
\begin{frame}[t]
  \frametitle{The Solution: Hong Kong's Regulatory Sandbox}
  A regulatory sandbox allows companies and public agencies to test selected LAE applications under controlled conditions.
  \begin{itemize}
    \item \textbf{Test safely:} Trials take place on approved routes and under specified operating conditions.
    \item \textbf{Learn early:} Companies and regulators collect evidence before wider deployment.
    \item \textbf{Reduce uncertainty:} Problems can be identified before very large investments are made.
    \item \textbf{Support regulation:} Trial results help the Government develop appropriate operating and safety requirements.
  \end{itemize}
  \textit{The sandbox does not remove safety requirements. It provides a controlled pathway for testing and learning.}
  \vfill
  {\tiny Source: Hong Kong Transport and Logistics Bureau, Low-altitude Economy Regulatory Sandbox X.}
\end{frame}

% ====================== Slide 21: STANDALONE FORMULA ======================
\begin{frame}[t]
  \frametitle{Applied Tool: How a Sandbox May Reduce Entry Barriers}
  \begin{block}{A Simple Framework}
  \begin{align*}
        \text{Entry Barrier}= & \text{Investment Cost}+\text{Compliance Cost}\\&+\text{Cost of Delay}+\text{Uncertainty}
  \end{align*}
  \end{block}
  \begin{itemize}
    \item \textbf{Investment cost}: Aircraft, infrastructure, software, and staff.
    \item \textbf{Compliance cost}: Testing, documentation, professional services, and approvals.
    \item \textbf{Cost of delay}: Cash spent while waiting before wider operation.
    \item \textbf{Uncertainty}: The risk that the technology, business model, or approval pathway may not work as expected.
  \end{itemize}
  \textbf{Core lesson:} A sandbox may reduce delay and uncertainty, making a project easier to test and finance. It does not eliminate compliance or safety costs.
\end{frame}

% ====================== Slide 22: BACK-OF-ENVELOPE (SANDBOX) ======================
\begin{frame}[t, shrink=10]
  \frametitle{Back-of-the-Envelope: The Value of Faster Learning}
  \textbf{Scenario:} A drone-delivery startup must decide how much capital to commit before it knows whether its service can operate reliably and gain approval.
  \begin{columns}[T]
    \begin{column}{0.48\textwidth}
      \begin{block}{Without Early Controlled Trials}
      \begin{itemize}
        \item Larger upfront commitment
        \item Slower operational learning
        \item More uncertainty about routes and requirements
        \item Greater risk of redesign after investment
      \end{itemize}
      \end{block}
    \end{column}
    \begin{column}{0.48\textwidth}
      \begin{block}{With a Regulatory Sandbox}
      \begin{itemize}
        \item Smaller staged tests
        \item Earlier evidence on safety and reliability
        \item Better information for regulators and investors
        \item Investment can be linked to milestones
      \end{itemize}
      \end{block}
    \end{column}
  \end{columns}
  \vspace{0.3cm}
  \textit{Takeaway: The main economic value of a sandbox is faster, lower-cost learning under controlled conditions.}
\end{frame}

% ====================== Slide 23: GBA COMPARATORS ======================
\begin{frame}[t]
  \frametitle{Working with the Greater Bay Area (GBA)}
  Hong Kong's LAE development is closely connected with the wider GBA ecosystem.
  \begin{itemize}
    \item \textbf{Manufacturing and technology:} Shenzhen and other GBA cities have strong capabilities in drones, electronics, batteries, and supply chains.
    \item \textbf{Hong Kong capabilities:} International finance, insurance, legal services, professional services, logistics, and a complex urban testing environment.
    \item \textbf{Cross-boundary opportunity:} Logistics and public-service applications may create value, but require coordination on routes, customs, data, communications, and safety requirements.
    \item \textbf{Strategic question:} How can Hong Kong and other GBA cities combine complementary strengths rather than duplicate them?
  \end{itemize}
\end{frame}

% ====================== Slide 24: CONCEPT PRIMER 3 ======================
\begin{frame}[t]
  \frametitle{Concept Primer: The "Empty Seat" Problem}
  
  \textbf{How do airlines make cheap tickets?}
  \begin{itemize}
      \item A Boeing 777 costs hundreds of millions of dollars. But it has \textbf{300+ seats}. By splitting the cost of the fuel and the pilots across 300 people, your ticket to Japan is affordable.
  \end{itemize}

  \vspace{0.3cm}
  \textbf{The Flying Taxi Problem:}
  \begin{itemize}
      \item An eVTOL flying taxi only has \textbf{4 seats}.
      \item You cannot divide the cost of the pilot and the machine by 300. You divide it by 4. 
      \item To prevent a ticket from costing HK\$5,000, you have to find a different mathematical trick.
  \end{itemize}
\end{frame}

% ====================== Slide 25: STANDALONE FORMULA ======================
\begin{frame}[t, shrink=10]
  \frametitle{Applied Tool: The Unit Cost Equation}
  \begin{block}{The Formula}
  \begin{align*}
         \text{Cost per Passenger}
    =& \text{Variable Cost per Passenger} \\
    +&\frac{\text{Daily Fixed Cost}}
    {\text{Flights per Day}\times\text{Seats per Flight}\times\text{Load Factor}}
  \end{align*}
  \end{block}
  \begin{itemize}
    \item \textbf{Utilization}: How frequently the aircraft flies.
    \item \textbf{Capacity}: How many passenger seats are available.
    \item \textbf{Load factor}: The percentage of seats actually occupied.
    \item \textbf{Core lesson}: Operators need both frequent operations and sufficient demand. Flying often does not help if the aircraft repeatedly flies with empty seats.
  \end{itemize}
\end{frame}

% ====================== Slide 26: BACK-OF-ENVELOPE (UNIT COST) ======================
\begin{frame}[t, shrink=15]
  \frametitle{Back-of-the-Envelope: Utilization and the Empty-Seat Problem}
  \textbf{Assumptions:} Fixed aircraft cost = HK\$10,000/day; variable cost = HK\$100/flight; four passenger seats.
  \begin{columns}[T]
    \begin{column}{0.48\textwidth}
      \begin{block}{Case A: 10 Flights, 50\% Load Factor}
      \begin{itemize}
        \item Paying passengers: \(10\times4\times50\%=20\)
        \item Fixed cost per passenger: HK\$500
        \item Variable cost per passenger: HK\$50
        \item \textbf{Total cost: HK\$550 per passenger}
      \end{itemize}
      \end{block}
    \end{column}
    \begin{column}{0.48\textwidth}
      \begin{block}{Case B: 25 Flights, 80\% Load Factor}
      \begin{itemize}
        \item Paying passengers: \(25\times4\times80\%=80\)
        \item Fixed cost per passenger: HK\$125
        \item Variable cost per passenger: HK\$31.25
        \item \textbf{Total cost: HK\$156.25 per passenger}
      \end{itemize}
      \end{block}
    \end{column}
  \end{columns}
  \textit{Takeaway: Technology may help, but demand, load factor, charging, maintenance, infrastructure, weather, and regulation all affect achievable utilization.}
  \vfill
  {\tiny Illustrative assumptions for classroom analysis; not current market data.}
\end{frame}

% ====================== Slide 27: TECH SUMMARY ======================
\begin{frame}[t]
  \frametitle{Technology: An Economic Tool}
  \begin{itemize}
    \item \textbf{Aircraft technology}: Changes speed, range, payload, capacity, reliability, and operating cost.
    \item \textbf{Automation}: May reduce crew requirements and support frequent operations, but must satisfy safety and regulatory requirements.
    \item \textbf{UTM and communications}: Support coordination, monitoring, and scalable operations in complex airspace.
    \item \textbf{Infrastructure}: Charging, maintenance, landing sites, and passenger processing determine actual throughput.
  \end{itemize}
  \textbf{Managerial lesson:} A technically impressive aircraft creates commercial value only when the complete operating system is safe, reliable, convenient, and affordable.
\end{frame}

% ====================== REVISED SLIDE 28: EXERCISE INSTRUCTIONS ======================
\begin{frame}[t]
  \frametitle{In-Class Exercise: Map the Money}
 
  \textbf{Activity: Map the Money in the Low Altitude Economy}
  \begin{itemize}
    \item \textbf{Format:} Groups of 4–5
    \item \textbf{Scenario:} Choose one:
      \begin{itemize}
        \item Hong Kong Airport Shuttle route
        \item Hospital Medical Delivery route
      \end{itemize}
    \item \textbf{Task:} On paper or in your notes, map the following:
    \begin{enumerate}
        \item Who are the key players in the LAE Value Chain?
        \item How does the money flow from the customer to each player?
        \item What does each player want? (Incentives and possible conflicts)
        \item How does Hong Kong’s Sandbox policy help this business succeed?
    \end{enumerate}
    \item \textbf{Timing:} 10 minutes group work + 5 minute quick share-out
  \end{itemize}
\end{frame}
% ====================== Slide 29: KEY TAKEAWAYS ======================
\begin{frame}[t]
  \frametitle{Key Takeaways}
  \begin{itemize}
    \item LAE combines aircraft, infrastructure, digital systems, operators, customers, and supporting services.
    \item A service creates value when it saves time, reduces cost or risk, improves access, or enables a mission that existing alternatives cannot perform well.
    \item Commercial viability depends on willingness to pay, demand, utilization, load factor, operating cost, and investment cost.
    \item Hong Kong may hold advantages in selected high-value services and in connecting capabilities across the GBA.
    \item A regulatory sandbox supports controlled testing and faster learning, but does not remove safety or compliance requirements.
    \item Technology is valuable when it improves the economics of the complete service, not merely the aircraft.
  \end{itemize}
\end{frame}

% ====================== Slide 30: ASSIGNMENT ======================
\begin{frame}[t]
  \frametitle{Week 1 Assignment Briefing}
  \textbf{Two-page Hong Kong LAE Economic Memo (maximum 800 words)}
  \begin{itemize}
    \item \textbf{Due:} Beginning of Week 2
    \item \textbf{Task:} Choose one Hong Kong LAE service idea, such as an airport shuttle, medical delivery, inspection service, or logistics route.
    \item \textbf{Requirements:}
    \begin{enumerate}
      \item Explain the customer problem and why someone would pay for the service.
      \item Estimate a plausible fare or service fee using time savings or another customer benefit.
      \item Estimate annual revenue: price $\times$ expected daily transactions $\times$ operating days.
      \item Identify two important barriers, such as demand, cost, infrastructure, regulation, safety, weather, or public acceptance.
      \item Explain how controlled sandbox testing could help address one barrier.
    \end{enumerate}
    \item \textbf{Grading focus:} Clear economic logic (40\%), transparent and realistic assumptions (30\%), policy link (20\%), professionalism (10\%).
  \end{itemize}
\end{frame}

% ====================== Slide 31: PREVIEW WEEK 2 ======================
\begin{frame}[t]
  \frametitle{Preview of Week 2}
  
  \begin{center}
      \Large \textbf{Market Structure: The Fight for the Sky}
  \end{center}
  
  \vspace{1cm}
  \textbf{Question to think about for next week:} \\
  \textit{If a landing pad on top of a Central skyscraper can only fit two drones at a time... who gets to land there, and who gets pushed out of business?}
\end{frame}

% ====================== Slide 32: THANK YOU ======================
\begin{frame}[t]
  \frametitle{Thank You + Q\&A}
  
  \textbf{Pre-Class Readings Reminder (for Week 2):}
  \begin{itemize}
      \item \textit{McKinsey Perspectives on Advanced Air Mobility}
      \item \textit{NASA UAM Market Study (2018)} 
  \end{itemize}
  
  \vspace{1cm}
  \begin{center}
      \textbf{Questions?} \\
      \small (Final 15 minutes reserved for extended Q\&A)
  \end{center}
\end{frame}

\end{document}